\documentclass[11pt]{article}
\usepackage{amsmath,amssymb,amsthm}
\usepackage[margin=1in]{geometry}
\usepackage{hyperref}

\newtheorem{theorem}{Theorem}[section]
\newtheorem{proposition}[theorem]{Proposition}
\newtheorem{corollary}[theorem]{Corollary}
\theoremstyle{definition}
\newtheorem{definition}{Definition}[section]
\theoremstyle{remark}
\newtheorem{remark}[theorem]{Remark}

\title{A Toy Paper on Even and Triangular Numbers}
\author{ProofLens example}
\date{}

\begin{document}
\maketitle

\begin{abstract}
A synthetic paper used to develop and test ProofLens. It contains a definition, a proved
theorem with hypotheses, a proved identity, a printed claim that is false together with its
corrected form, a theorem with a uniform constant, a claim left open, and a remark outside the
formal scope.
\end{abstract}

\section{Definitions}

\begin{definition}[Even numbers]\label{def:even}
A natural number $n$ is \emph{even} if $n = 2k$ for some natural number $k$. The
\emph{triangular number} $T_n$ is $T_n = 0 + 1 + \dots + n$.
\end{definition}

\section{Results}

\begin{theorem}[Closure under addition]\label{thm:even-add}
If $a$ and $b$ are even natural numbers, then $a + b$ is even.
\end{theorem}

\begin{proof}
Write $a = 2j$ and $b = 2k$; then $a + b = 2(j + k)$.
\end{proof}

\begin{theorem}[Gauss]\label{thm:triangle}
For every natural number $n$,
\begin{equation}\label{eq:gauss}
  2\,T_n = n(n+1).
\end{equation}
\end{theorem}

\begin{proof}
Induction on $n$: $2T_{n+1} = 2T_n + 2(n+1) = n(n+1) + 2(n+1) = (n+1)(n+2)$.
\end{proof}

\begin{proposition}\label{prop:pred}
Every natural number is strictly greater than its predecessor $n - 1$.
\end{proposition}

\begin{remark}\label{rem:pred}
With truncated subtraction $0 - 1 = 0$, Proposition~\ref{prop:pred} fails at $n = 0$; it holds
for every $n \ge 1$.
\end{remark}

\begin{theorem}[Uniform bound]\label{thm:bound}
There is a constant $C$, independent of $n$, such that $T_n \le C n^2$ for every natural number
$n$.
\end{theorem}

\begin{proof}
By Theorem~\ref{thm:triangle}, $2T_n = n^2 + n \le 2n^2$, so $C = 1$ works.
\end{proof}

\begin{corollary}\label{cor:even-triangle}
For every natural number $n$, the triangular number $T_{4n}$ is even.
\end{corollary}

\begin{proof}
$T_{4n} = 2n(4n + 1)$ by \eqref{eq:gauss}.
\end{proof}

\section{Discussion}

\begin{remark}\label{rem:history}
The identity of Theorem~\ref{thm:triangle} is traditionally attributed to Gauss.
\end{remark}

\end{document}
